\chapter{Conclusion and Future Work}
\label{ch:conclusion}

\section{Summary}

This thesis set out to make charts carry verified statistical findings, and to test
whether a model can learn to state those findings from the image alone. We built
InsightVLM-20K, 20{,}000 charts in 26 chart types whose labels are issued only when a
named statistical test confirms them on exactly the plotted data, with 6{,}083 charts
from 654 real public tables held in source-disjoint splits and a supplementary pool of
2{,}545 further real training charts.

\begin{description}[leftmargin=1.2cm,style=nextline,font=\bfseries]
  \item[RQ1.] A fine-tuned CNN ensemble identifies the finding of a real chart with
  0.867 accuracy (95\% CI $[0.846, 0.886]$) and of a synthetic chart with 0.999, without
  a supplied question. Zero-shot vision--language models reach at most 0.468.
  \item[RQ2.] A deterministic pipeline on the data table is right 99.95\% of the time and
  exact. A QLoRA-tuned generative VLM names the finding 78.7\% of the time but cites
  correct numbers only 13.9\% of the time.
  \item[RQ3.] The model's Grad-CAM peak falls inside a verified evidence box on 52.4\% of real
  charts against 43.9\% for a random point: grounded beyond chance for time series, group
  comparisons and relationships, but not for distributions.
  \item[RQ4.] Real data costs 13 points relative to synthetic data. A plotting library
  never seen in training costs a further 15 points (0.716), reduced but not removed by
  training on charts drawn by more than one library.
  \item[RQ5.] Accuracy should be reported on real charts, across more than one renderer,
  against a pixel-free floor (0.288) and a data-oracle ceiling (0.9995), with selection
  on validation only; and trained models must be tested for shortcuts directly, because
  a data audit missed the one that mattered.
\end{description}

Along the way we found and removed a rendering leak that had inflated accuracy on
real time-series charts to 0.94--0.96 from a true 0.31--0.41, and corrected an earlier result that had been
selected on the test set. The final numbers are lower than the earliest drafts reported,
and they are the ones we can defend.

\section{Deliverables}

\begin{itemize}[leftmargin=1.3em]
  \item InsightVLM-20K version 2, with datasheet, bias audit, the supplementary real pool,
  and Plotly and Vega-Lite re-renders of the real test charts.
  \item Trained models and the deployed ensemble.
  \item An interactive tester (Appendix~\ref{app:tester}) that runs uploaded charts
  through the evaluated pipeline.
  \item The full build, training and evaluation code, from which every number in this
  thesis is regenerated.
\end{itemize}

\section{Future Work}

\begin{enumerate}[leftmargin=1.4em]
  \item \textbf{Extend the human study} to more independent annotators, add the human
  classification baseline on 200 charts, and test whether labelling the subtle classes by
  visual salience as well as by statistics changes what models learn.
  \item \textbf{Close the unseen-renderer gap} with more rendering libraries in training
  (Excel-style, ggplot2, D3) and with charts taken from real documents.
  \item \textbf{Scale the generative model} to a larger backbone and the full training set,
  and optimise it for numeric fidelity, not only class accuracy.
  \item \textbf{Multi-label findings}, so a chart can show a trend and an anomaly at once.
  \item \textbf{More real data for rare classes} (heavy tails, seasonality, left skew,
  inverted-U relationships), which are also the weakest classes.
  \item \textbf{Seed variance for every configuration}, not only the deployed recipe, so
  ablation differences can be judged against training variability.
\end{enumerate}
